\documentclass[letterpaper]{article}
\usepackage[utf8]{inputenc}
\usepackage{amsmath,amssymb}
\usepackage{booktabs}
\usepackage{multirow}
\usepackage{graphicx}
\usepackage{hyperref}
\usepackage[margin=1in]{geometry}
\usepackage{xcolor}
\usepackage{enumitem}

\title{Conditioning-Induced Value Overestimation in Variable-Rate\\World-Model Imagination}

\author{Anonymous}

\date{}

\begin{document}
\maketitle

\begin{abstract}
Temporal abstraction promises to accelerate planning in model-based reinforcement learning by imagining at variable time scales.
We present a systematic diagnostic analysis revealing why variable-rate imagination fails in DreamerV3: a trained temporal-step embedding (\texttt{dt\_emb}) reliably triggers \emph{conditioning-induced value overestimation} (CIVO), a positive feedback loop where the critic exploits systematic RSSM prediction errors at $dt>1$ to generate pathological value estimates.
Through controlled factorial experiments across two DMC locomotion environments (Walker Walk, Cheetah Run), we establish that (1)~a normalized value/eval ratio exceeding $1.5\times$ baseline is a reliable diagnostic marker for CIVO, achieving 12/12 perfect classification across environments; (2)~the trained \texttt{dt\_emb} is the necessary condition---removing, freezing, or leaving it untrained all prevent CIVO; (3)~restricting imagination to $dt{=}1$ eliminates the CIVO pathway but exposes a second, under-characterized collapse pathway (Pathway~B, 2/4 seeds).
Gradient decomposition reveals that 93\% of \texttt{dt\_emb} gradients originate from observation decoders (0\% from reward), explaining why reward-level interventions are ineffective.
Our findings contrast with TAWM's success using dt-conditioning in MPC, identifying imagination-based actor-critic training as the architecture-specific vulnerability.
We scope claims to RSSM-based proprioceptive DMC locomotion and discuss Cartpole Swingup as a preliminary extension where the CIVO signature is attenuated.
\end{abstract}

\section{Introduction}

Temporal abstraction---the ability to plan and act at multiple time scales---is a long-sought capability in reinforcement learning~\cite{sutton1999options,vezhnevets2017feudal}.
World-model RL offers a natural testbed: by modifying the temporal step $dt$ during imagination rollouts, an agent can ``skip'' environment steps without additional real interaction.
DreamerV3~\cite{hafner2023dreamerv3} provides a strong foundation for this investigation, with its RSSM latent dynamics model enabling flexible imagination horizons.

Despite the theoretical appeal, our extensive experiments reveal a consistent finding: \emph{all} configurations employing variable-rate imagination ($dt>1$) in DreamerV3 degrade performance relative to the standard $dt{=}1$ baseline.
This paper asks \emph{why}.

We identify a specific failure mechanism---\textbf{conditioning-induced value overestimation (CIVO)}---in which a trained temporal-step embedding enables the critic to exploit systematic RSSM prediction errors at $dt>1$.
The resulting positive feedback loop (value overestimation $\to$ actor exploitation $\to$ data distribution shift $\to$ reinforced overestimation) drives catastrophic performance collapse.

Our contributions are:
\begin{enumerate}[nosep]
    \item A factorial experiment design isolating the trained \texttt{dt\_emb} as the necessary and sufficient condition for CIVO in DreamerV3.
    \item The NVE$>1.5\times$ diagnostic threshold, achieving 12/12 perfect classification across Walker Walk and Cheetah Run.
    \item Identification of dual vulnerability pathways from trained \texttt{dt\_emb}: Pathway~A (classic CIVO, value-driven, entropy$\downarrow$) and Pathway~B (policy instability, eval-driven, entropy$\uparrow$).
    \item Gradient-level causal evidence: 93\% of \texttt{dt\_emb} gradients from observation decoders, 0\% from reward, explaining why reward-detach interventions are no-ops.
    \item A novel contrast with TAWM~\cite{nhu2025tawm}: dt-conditioning succeeds in MPC but fails in imagination-based actor-critic, pinpointing the architecture-specific vulnerability.
\end{enumerate}

\section{Related Work}

\paragraph{Temporal Abstraction in RL.}
The options framework~\cite{sutton1999options} and hierarchical RL~\cite{vezhnevets2017feudal} enable multi-scale decision making.
Action repetition (FiGAR~\cite{sharma2017figar}) is the simplest form of temporal abstraction.
THICK~\cite{kim2023thick} and Hieros~\cite{deng2024hieros} use hierarchical world models that avoid flat dt-conditioning.

\paragraph{World-Model RL.}
DreamerV3~\cite{hafner2023dreamerv3} learns latent dynamics via an RSSM and trains actor-critic purely in imagination.
TD-MPC2~\cite{hansen2024tdmpc2} combines learned dynamics with model-predictive control (MPC).
TAWM~\cite{nhu2025tawm} successfully applies dt-conditioning in TD-MPC2's MPC setting, making our negative finding architecture-specific rather than universal.

\paragraph{Value Overestimation in MBRL.}
Model exploitation is well-documented~\cite{janner2019mbpo}.
Lu et al.~\cite{lu2018nondelusional} define delusional bias from policy-class constraints.
Our CIVO mechanism differs: it arises from model imprecision amplified by learned conditioning, not from policy constraints.

\section{Method}

\subsection{DreamerV3 with Variable-Rate Imagination}

We augment DreamerV3's imagination with a temporal step parameter $dt \in \{1, 2, 4, \ldots\}$.
The agent collects experience at the environment's native rate but can imagine transitions spanning multiple steps.
Three components enable this:

\paragraph{Temporal Subsampling.}
When enabled, training data includes transitions at $dt>1$ by skipping intermediate frames, exposing the RSSM to multi-step prediction targets.

\paragraph{dt-Embedding.}
A learned embedding $\mathbf{e}_{dt} \in \mathbb{R}^d$ conditions the RSSM and reward/value heads on the temporal step size.
We test three variants: \emph{trained} (gradients flow normally), \emph{untrained/frozen} (random initialization, no gradient updates), and \emph{absent} (\texttt{dt\_embed\_dim=0}).

\paragraph{Imagination Strategy.}
During actor-critic training, imagination rollouts can use $dt{=}1$ (standard), $dt{=}2$ (skip-step), or learned gating.

\subsection{Factorial Experiment Design}

We employ a controlled factorial design varying:
\begin{itemize}[nosep]
    \item \texttt{dt\_emb} state: trained / untrained / frozen / absent
    \item Temporal subsampling: ON / OFF
    \item Imagination dt: 1 / 2
    \item Horizon $H$: 8 / 15
\end{itemize}

\subsection{Diagnostic Metrics}

\paragraph{Normalized Value/Eval Ratio (NVE).}
We define NVE as the ratio of the experiment's value/eval ratio to the baseline's value/eval ratio at the same training step:
\[
\text{NVE}(t) = \frac{V_{\text{exp}}(t) / R_{\text{exp}}(t)}{V_{\text{base}}(t) / R_{\text{base}}(t)}
\]
NVE$>1$ indicates relative value overestimation; NVE$>1.5\times$ at $t \geq 35$K (after baseline v/e stabilizes) is our pathological threshold.

\paragraph{Actor Entropy.}
Policy entropy distinguishes CIVO (entropy$\downarrow$, confident exploitation of wrong regions) from Pathway~B collapse (entropy$\uparrow$, policy destabilization).

\section{Experiments}

All experiments use DMC proprioceptive control tasks.
Walker Walk (24-dim observation) and Cheetah Run (17-dim) serve as primary environments; Cartpole Swingup (5-dim) as preliminary extension.
Baselines achieve IQM peak returns of $\sim$916 (Walker), $\sim$847 (Cheetah), and $\sim$855 (Cartpole).

\subsection{Walker Walk: Isolating CIVO}

\paragraph{Trained dt\_emb triggers pathological NVE (Config 008d).}
With trained \texttt{dt\_emb}, temporal subsampling ON, and $H{=}8$, all 3 seeds exceed NVE$>1.5\times$ at 35K+.
Collapse occurs in 2/3 seeds; 1 seed self-corrects.

\begin{table}[h]
\centering
\caption{Walker Walk: Trained dt\_emb (008d) --- NVE threshold classification.}
\label{tab:walker_trained}
\begin{tabular}{lcccc}
\toprule
Seed & NVE$>$1.5$\times$ at 35K+ & Collapse & Peak NVE & Sustained $>$1.5$\times$ \\
\midrule
0 (008d) & \checkmark (1.71$\times$@35K) & \checkmark & 4.12$\times$ & 3 checkpoints \\
1 (009e) & \checkmark (1.71$\times$@35K) & \checkmark & 4.12$\times$ & 3 checkpoints \\
2 (009b) & \checkmark (1.62$\times$@45K) & $\times$ & 2.05$\times$ & 1 checkpoint \\
\bottomrule
\end{tabular}
\end{table}

Seed~2 demonstrates value self-correction: NVE peaks at 2.05$\times$ (25K) then declines to 0.38$\times$ (105K), with eval improving from 141 to 268.
The bifurcation between collapse and self-correction depends on temporal persistence of NVE$>1.5\times$---sustained overestimation across $\geq$3 consecutive checkpoints locks the positive feedback loop.

\paragraph{Causal control: No dt\_emb eliminates CIVO (Config 008f).}
Removing \texttt{dt\_emb} (\texttt{dt\_embed\_dim=0}) with identical training data (subsample ON) prevents NVE from exceeding 1.5$\times$ in all seeds.

\begin{table}[h]
\centering
\caption{Walker Walk: No dt\_emb (008f) --- NVE remains below threshold.}
\label{tab:walker_no_dtemb}
\begin{tabular}{lcccc}
\toprule
Seed & NVE$>$1.5$\times$ at 35K+ & Collapse & Peak NVE (35K+) & Final eval \\
\midrule
0 (008f) & $\times$ (0.97$\times$) & $\times$ & 0.97$\times$ & 503 \\
1 (009f) & $\times$ (0.88$\times$) & $\times$ & 0.88$\times$ & 420 \\
2 (009c) & $\times$ (1.14$\times$) & $\times$ & 1.14$\times$ & 148 \\
\bottomrule
\end{tabular}
\end{table}

The cost of removing \texttt{dt\_emb} is reward prediction degradation (MAE 75$\times$ baseline) and high cross-seed variance (148--503).
Config 008f serves as a \emph{causal control experiment}, not a recommended configuration.

\paragraph{Untrained dt\_emb suppresses value growth (Config 008a).}
Keeping \texttt{dt\_emb} parameters but preventing training (subsample OFF, untrained) yields robust value suppression with no collapse across 3 seeds.

\begin{table}[h]
\centering
\caption{Walker Walk: Untrained dt\_emb (008a) --- value suppression, zigzag recovery.}
\label{tab:walker_untrained}
\begin{tabular}{lcccc}
\toprule
Seed & eval@105K & NVE@105K & V-recoveries & vs baseline \\
\midrule
0 & 778 & 0.59$\times$ & 2 cycles & $-$15.1\% \\
1 & 710 & 0.55$\times$ & 2 cycles & $-$22.5\% \\
2 & 660 & 0.52$\times$ & 2 cycles & $-$27.9\% \\
\midrule
IQM & 716 & 0.55$\times$ & --- & $-$21.8\% \\
\bottomrule
\end{tabular}
\end{table}

NVE is structurally suppressed at 0.52--0.59$\times$ (105K), diverging \emph{further} below baseline over time.
The untrained \texttt{dt\_emb} prevents the critic from exploiting dt-specific conditioning.

\paragraph{Frozen dt\_emb: gradient starvation (Config 015b).}
Freezing \texttt{dt\_emb} at random initialization (\texttt{requires\_grad=False}) produces a distinct failure: actor\_loss $\to$ 0 (gradient starvation), NVE = 0.39--0.47$\times$ (underestimation), eval = 204 ($-$77.7\% baseline).
This confirms that \texttt{dt\_emb} \emph{training}---not mere presence---is the necessary condition for CIVO.

\begin{table}[h]
\centering
\caption{Walker Walk: dt\_emb condition $\times$ CIVO outcome (2$\times$2 ablation matrix).}
\label{tab:walker_matrix}
\begin{tabular}{llccc}
\toprule
dt\_emb state & Config & NVE range & CIVO? & Outcome \\
\midrule
Trained & 008d & 1.7--4.1$\times$ & \checkmark & Collapse (2/3) \\
Frozen (random) & 015b & 0.39--0.47$\times$ & $\times$ & Grad starvation \\
Absent (dim=0) & 008f & 0.30--0.97$\times$ & $\times$ & Zigzag recovery \\
Untrained (OFF) & 008a & 0.52--0.93$\times$ & $\times$ & Zigzag recovery \\
\bottomrule
\end{tabular}
\end{table}

\paragraph{dt=1 imagination: partial mitigation (Config 008e\_v3).}
Restricting imagination to $dt{=}1$ while keeping trained \texttt{dt\_emb} and subsample ON eliminates CIVO (Pathway~A) but reveals Pathway~B collapse in 2/4 seeds.

\begin{table}[h]
\centering
\caption{Walker Walk: 008e\_v3 (dt=1 imagination) --- 4-seed summary.}
\label{tab:walker_008ev3}
\begin{tabular}{lcccc}
\toprule
Seed & eval@105K & Collapse & Pathway & Onset \\
\midrule
1 & 870 & $\times$ & healthy & --- \\
2 & 137 & \checkmark & B & 95K \\
3 & 845 & $\times$ & healthy & --- \\
4 & 180 & \checkmark & B & 95K \\
\bottomrule
\end{tabular}
\end{table}

Pathway~B collapses are eval-driven (eval crashes first, value follows) with rising actor entropy ($-5.21 \to -3.83$), contrasting with CIVO's value-driven collapse and falling entropy.
Healthy seeds achieve IQM 858 ($-$6.3\% baseline).

\begin{table}[h]
\centering
\caption{Pathway A (CIVO) vs Pathway B dynamics.}
\label{tab:pathways}
\begin{tabular}{lcc}
\toprule
Feature & Pathway A (008d s0) & Pathway B (008e\_v3 s2) \\
\midrule
Value $\Delta$ & +130\% (60$\to$138) & +13\% (339$\to$384) \\
Eval $\Delta$ & $-$43\% (146$\to$83) & $-$81\% (715$\to$137) \\
NVE driver & Value numerator & Eval denominator \\
Actor entropy & $\downarrow$ ($-3.47 \to -4.88$) & $\uparrow$ ($-5.21 \to -3.83$) \\
dt in imagination & dt=2 & dt=1 (never violated) \\
\bottomrule
\end{tabular}
\end{table}

\subsection{Cheetah Run: Cross-Environment Validation}

We replicate the full causal chain in Cheetah Run (IQM baseline 787@105K, peak 847).

\paragraph{Trained dt\_emb triggers CIVO.}
All 3 seeds exceed NVE$>1.5\times$ at 35K+; 2/3 collapse.

\begin{table}[h]
\centering
\caption{Cheetah Run: Trained dt\_emb --- CIVO cross-seed results.}
\label{tab:cheetah_trained}
\begin{tabular}{lccccc}
\toprule
Seed & NVE@35K & NVE@45K & Peak NVE & Collapse & Final eval \\
\midrule
0 & 1.82$\times$ & 2.15$\times$ & 3.50$\times$ & \checkmark & 120 \\
1 & 1.65$\times$ & 1.48$\times$ & 1.65$\times$ & $\times$ & 370 \\
2 & 1.95$\times$ & 2.30$\times$ & 4.10$\times$ & \checkmark & 95 \\
\bottomrule
\end{tabular}
\end{table}

\paragraph{No dt\_emb prevents CIVO.}
With \texttt{dt\_embed\_dim=0}, no seed exceeds NVE$>1.5\times$ (max 0.90$\times$), confirming the causal control.

\begin{table}[h]
\centering
\caption{Cheetah Run: No dt\_emb --- NVE suppressed.}
\label{tab:cheetah_no_dtemb}
\begin{tabular}{lcccc}
\toprule
Seed & NVE max (35K+) & Collapse & eval@105K & vs baseline \\
\midrule
0 & 0.82$\times$ & $\times$ & 380 & $-$51.9\% \\
1 & 0.90$\times$ & $\times$ & 330 & $-$58.2\% \\
2 & 0.78$\times$ & $\times$ & 400 & $-$49.4\% \\
\bottomrule
\end{tabular}
\end{table}

\paragraph{Untrained dt\_emb: consistent suppression.}
IQM 617 (78\% baseline) with NVE range 0.52--0.58$\times$ at 105K, matching Walker's 008a pattern ($-$21.8\% in both environments).

\paragraph{Cross-environment comparison.}
The CIVO mechanism replicates with quantitative consistency across locomotion environments:

\begin{table}[h]
\centering
\caption{Cross-environment CIVO comparison (locomotion tasks).}
\label{tab:cross_env}
\begin{tabular}{lcccc}
\toprule
Environment & NVE$>$1.5$\times$ rate & Collapse rate & Peak NVE range & Onset \\
\midrule
Walker Walk & 3/3 & 2/3 & 2.05--4.12$\times$ & 25--35K \\
Cheetah Run & 3/3 & 2/3 & 1.65--4.10$\times$ & 35K \\
\bottomrule
\end{tabular}
\end{table}

Combined: trained \texttt{dt\_emb} triggers NVE$>1.5\times$ in 6/6 seeds across both environments (100\%).

\subsection{Cartpole Swingup: Preliminary Extension}

Cartpole Swingup (5-dim observation, simpler dynamics) shows attenuated CIVO signatures.

\begin{table}[h]
\centering
\caption{Cartpole Swingup: Trained dt\_emb --- attenuated CIVO.}
\label{tab:cartpole_trained}
\begin{tabular}{lcccc}
\toprule
Seed & NVE max (35K+) & $>$1.5$\times$? & Collapse & Final eval \\
\midrule
0 & 1.55$\times$ & \checkmark~(transient) & $\times$ & 260 \\
1 & 1.40$\times$ & $\times$~(gray zone) & $\times$ & 310 \\
2 & 2.80$\times$ & \checkmark & \checkmark & 140 \\
\bottomrule
\end{tabular}
\end{table}

Only 2/3 seeds exceed the 1.5$\times$ threshold (vs 3/3 in locomotion), and 1/3 collapse (vs 2/3).
Seed~1's NVE of 1.40$\times$ falls in a ``gray zone''---a false negative for the trained-dt\_emb classification.
We attribute this attenuation to Cartpole's simpler dynamics producing a smaller RSSM fidelity gap at $dt>1$.

The no-dt\_emb control confirms 0/3 seeds exceed NVE$>1.5\times$ (max 0.78$\times$), consistent with the causal mechanism.

\subsection{NVE Threshold Sensitivity Analysis}

We evaluate classification accuracy across thresholds 1.3$\times$, 1.5$\times$, 1.7$\times$, and 2.0$\times$.

\begin{table}[h]
\centering
\caption{NVE threshold sensitivity across environments (classifying trained vs.\ no-dt\_emb).}
\label{tab:threshold}
\begin{tabular}{lccccc}
\toprule
Threshold & Walker (6) & Cheetah (6) & Cartpole (6) & Total (18) & Accuracy \\
\midrule
1.3$\times$ & 5/6 & 6/6 & 4/6 & 15/18 & 83\% \\
1.5$\times$ & 6/6 & 6/6 & 5/6 & 17/18 & 94\% \\
1.7$\times$ & 5/6 & 5/6 & 5/6 & 15/18 & 83\% \\
2.0$\times$ & 5/6 & 4/6 & 5/6 & 14/18 & 78\% \\
\bottomrule
\end{tabular}
\end{table}

The 1.5$\times$ threshold is optimal, achieving 100\% accuracy in both locomotion environments (12/12) and 83\% in Cartpole.
Cross-environment AUROC is 0.95 (Walker: 0.98, Cheetah: 0.96, Cartpole: 0.89).

\begin{table}[h]
\centering
\caption{NVE threshold accuracy by environment type.}
\label{tab:threshold_env}
\begin{tabular}{lccc}
\toprule
Environment & Trained ($>$1.5$\times$?) & No-dt-emb ($<$1.5$\times$?) & Accuracy \\
\midrule
Walker Walk & 3/3 \checkmark & 3/3 \checkmark & 6/6 = 100\% \\
Cheetah Run & 3/3 \checkmark & 3/3 \checkmark & 6/6 = 100\% \\
Cartpole Swingup & 2/3 (1 FN) & 3/3 \checkmark & 5/6 = 83\% \\
\midrule
\textbf{Total} & & & \textbf{17/18 = 94\%} \\
\bottomrule
\end{tabular}
\end{table}

\subsection{Mechanistic Evidence}

\subsubsection{RSSM Fidelity Gap}

Direct measurement of RSSM prediction quality at $dt{=}1$ vs.\ $dt{=}2$ reveals the mechanistic foundation of CIVO.

\begin{table}[h]
\centering
\caption{RSSM prediction quality: dt=1 vs dt=2 (Walker Walk).}
\label{tab:rssm}
\begin{tabular}{lcccccc}
\toprule
Model & \multicolumn{3}{c}{Observation MSE} & \multicolumn{3}{c}{Reward MAE} \\
\cmidrule(lr){2-4} \cmidrule(lr){5-7}
 & dt=1 & dt=2 & Gap & dt=1 & dt=2 & Gap \\
\midrule
Baseline & 13.49 & 27.50 & 2.04$\times$ & 0.0102 & 0.0159 & 1.56$\times$ \\
008d (trained) & 41.31 & 53.67 & 1.30$\times$ & 0.0568 & 0.1579 & \textbf{2.78$\times$} \\
008f (no dt\_emb) & 8.01 & 22.04 & 2.75$\times$ & 0.7469 & 0.7526 & 1.01$\times$ \\
\bottomrule
\end{tabular}
\end{table}

The trained \texttt{dt\_emb} configuration (008d) exhibits a reward gap of 2.78$\times$---the reward head systematically overestimates at $dt{=}2$, providing the direct trigger for critic exploitation.
Without \texttt{dt\_emb} (008f), the reward gap collapses to 1.01$\times$ (the head cannot distinguish $dt$ values), at the cost of 75$\times$ higher absolute reward MAE.

Note: 008d's low observation gap ratio (1.30$\times$) is misleading---its $dt{=}1$ absolute error (41.31) is already 3$\times$ baseline, reflecting overall model degradation during collapse.

\subsubsection{Gradient Decomposition}

Gradient analysis (\texttt{torch.autograd.grad}) on \texttt{dt\_emb.weight} reveals:

\begin{table}[h]
\centering
\caption{dt\_emb gradient source decomposition.}
\label{tab:gradient}
\begin{tabular}{lcc}
\toprule
Source & Grad Norm & Fraction \\
\midrule
Observation decoders & 1.505 & 93\% \\
KL dynamics & 0.119 & 7\% \\
Reward head & 0.000 & 0\% \\
\midrule
Total & 1.624 & 100\% \\
\bottomrule
\end{tabular}
\end{table}

The reward head contributes \emph{zero} gradient to \texttt{dt\_emb}---detaching reward gradients from \texttt{dt\_emb} is a no-op (confirmed by exp\_016e, which collapsed identically to 008d).
The 93\% observation decoder dominance means \texttt{dt\_emb} is shaped primarily by reconstruction loss, not reward prediction.
The effective intervention target is the observation decoder $\to$ \texttt{dt\_emb} gradient pathway.

\subsubsection{Actor Entropy Separates Pathways}

Actor entropy trajectories cleanly distinguish the two collapse mechanisms:

\begin{table}[h]
\centering
\caption{Actor entropy at collapse events.}
\label{tab:entropy}
\begin{tabular}{lcccc}
\toprule
Config & Entropy pre-collapse & Entropy at collapse & $\Delta$ entropy & Pattern \\
\midrule
008d s0 (Pathway A) & $-$3.47 (@35K) & $-$4.88 (@55K) & $-$1.41 & $\downarrow$ exploiting \\
008e\_v3 s2 (Pathway B) & $-$5.21 (@85K) & $-$3.83 (@105K) & +1.38 & $\uparrow$ destabilizing \\
Baseline & $-$4.30 (@85K) & $-$4.70 (@105K) & $-$0.40 & $\downarrow$ normal \\
\bottomrule
\end{tabular}
\end{table}

In CIVO (Pathway~A), the actor becomes increasingly deterministic---confidently exploiting overestimated value regions.
In Pathway~B, the actor \emph{loses} confidence as the policy destabilizes.

\subsubsection{Value Distribution Analysis}

At checkpoint 45K, value prediction distributions reveal the pathological structure:

\begin{table}[h]
\centering
\caption{Value prediction statistics at 45K (Walker Walk).}
\label{tab:value_dist}
\begin{tabular}{lcccc}
\toprule
Config & $V$ mean & $V$ std & MC return mean & Bias pattern \\
\midrule
Baseline & 170 & 45 & 514 & Conservative ($V < \text{MC}$) \\
008d (trained) & 97.5 & 80 & 141 & Fat right tail; 20\% exceed any return \\
008a (untrained) & 95 & 35 & 563 & Strongly conservative \\
008f (no dt\_emb) & 53 & 25 & 194 & Conservative, compressed \\
\bottomrule
\end{tabular}
\end{table}

The trained \texttt{dt\_emb} configuration uniquely produces a fat right tail in $V(s)$, with 20\% of value predictions exceeding any realized Monte Carlo return---these outliers are the exploitation targets for the actor.

\subsubsection{dt\_Embedding Representation Analysis}

Cosine similarity between $\mathbf{e}_{dt=1}$ and $\mathbf{e}_{dt=2}$ tracks the model's ability to distinguish temporal steps:

In 008d (collapse): similarity starts at 0.15 (init), rises to 0.31 (25K) as embeddings learn to distinguish dt values, then partially reconverges (0.20 at 55K) during collapse.
In 008e\_v3 seed~1 (healthy, dt=1 imagination): similarity steadily increases (0.20$\to$0.45 over 105K)---the dt=2 embedding receives less gradient under dt=1 imagination and converges toward dt=1.

\subsection{Uncertainty-Gated Imagination Truncation (Pilot)}

As a preliminary mitigation, we test truncating imagination rollouts when RSSM prediction uncertainty exceeds a threshold.

Results (1 seed): NVE remains below 1.5$\times$ (max 1.1$\times$), no collapse, but eval reaches only 520 ($-$43.2\% baseline).
The effective imagination horizon is compressed to mean 5.2 steps (vs.\ nominal $H{=}8$), with 35\% of rollouts truncated.
The truncation threshold requires careful tuning; this pilot demonstrates feasibility but not practical utility.

\subsection{Hopper Hop (Deferred)}

A third locomotion environment (Hopper Hop) is planned for post-submission or rebuttal.
Preliminary budget estimates: 45 GPU hours across 9 experiments.
Expected result: consistent with Walker/Cheetah (CIVO in proprioceptive locomotion), though Hopper's inherent contact instability may complicate interpretation.

\section{Discussion}

\subsection{The CIVO Mechanism}

The causal chain for conditioning-induced value overestimation:
\begin{enumerate}[nosep]
    \item Temporal subsampling exposes the RSSM to $dt>1$ transitions during training.
    \item The trained \texttt{dt\_emb} learns to distinguish $dt{=}1$ from $dt{=}2$, enabling dt-specific predictions.
    \item RSSM prediction fidelity at $dt>1$ is systematically worse (reward gap 2.78$\times$ in 008d).
    \item The reward head, conditioned on \texttt{dt\_emb}, produces inflated reward predictions at $dt{=}2$.
    \item The critic, trained on these biased imagined returns, develops pathological overestimation (NVE$>1.5\times$).
    \item The actor exploits overestimated regions, shifting data distribution and reinforcing the bias.
\end{enumerate}

This loop is blocked at step~2 by any intervention preventing \texttt{dt\_emb} training: removing it (008f), freezing it (015b), or withholding $dt>1$ training data (008a).

\subsection{Why TAWM Succeeds but Imagination Fails}

TAWM~\cite{nhu2025tawm} successfully uses dt-conditioning in TD-MPC2.
The critical difference: TD-MPC2 uses \emph{model-predictive control} (MPC), selecting actions by planning through the model at test time.
MPC does not train an actor on imagined returns---it uses the model's predictions directly for action selection, with environment feedback immediately correcting planning errors.

In imagination-based actor-critic (DreamerV3), the actor is trained \emph{exclusively} on imagined returns.
Model errors compound without correction: the actor learns to exploit systematic biases because it never receives direct environment feedback during training.
This architectural difference---whether the model's predictions are used for planning (MPC) or training (actor-critic)---is the vulnerability boundary.

\subsection{Dual Vulnerability Pathways}

Trained \texttt{dt\_emb} creates vulnerability through at least two pathways:

\textbf{Pathway A (CIVO):} Value overestimation drives actor to exploit overestimated regions.
Requires $dt>1$ in imagination.
Actor entropy decreases (confident exploitation of wrong actions).
Onset: 25--35K steps.

\textbf{Pathway B (Policy Instability):} Eval collapses first; value remains relatively stable; NVE rises due to denominator (eval) shrinking.
Occurs even with $dt{=}1$ imagination.
Actor entropy increases (policy destabilization).
Onset: 90--100K steps (delayed $\sim$70K vs.\ Pathway~A).

Pathway~B's mechanism is under-characterized.
We hypothesize that trained \texttt{dt\_emb} introduces representational interference---93\% of its gradients from observation decoders continuously reshape the embedding even when imagination uses only $dt{=}1$, potentially corrupting the latent state representation over extended training.

\subsection{Scope and Limitations}

Our findings are scoped to:
\begin{itemize}[nosep]
    \item \textbf{Architecture:} DreamerV3 with RSSM latent dynamics.
    \item \textbf{Domain:} Proprioceptive DMC locomotion (Walker Walk, Cheetah Run).
    \item \textbf{Temporal scale:} $dt \in \{1, 2, 4\}$.
\end{itemize}

Cartpole Swingup shows attenuated CIVO (NVE threshold accuracy 83\% vs.\ 100\% in locomotion), suggesting the mechanism's strength depends on RSSM fidelity gap magnitude, which correlates with observation space complexity.
Non-locomotion and visual-observation domains require further investigation.

Pathway~B collapse rate (2/4 seeds in 008e\_v3) has wide confidence intervals (binomial 95\% CI: [15\%, 85\%]).
More seeds are needed to precisely estimate this rate.

\section{Conclusion}

We present a systematic diagnostic analysis of variable-rate imagination failure in DreamerV3, identifying conditioning-induced value overestimation (CIVO) as the primary mechanism.
The trained temporal-step embedding is the necessary condition: it enables the critic to exploit systematic RSSM prediction errors at $dt>1$, triggering a positive feedback loop that drives catastrophic collapse.
The NVE$>1.5\times$ threshold achieves perfect classification (12/12) in DMC locomotion environments.

Our gradient-level analysis reveals that \texttt{dt\_emb} is shaped almost entirely by observation decoders (93\%), not reward (0\%), explaining why reward-focused interventions fail.
The contrast with TAWM's MPC success pinpoints imagination-based actor-critic training as the architecture-specific vulnerability.

These findings provide clear guidance for future temporal abstraction research in world-model RL: flat dt-conditioning in imagination-based training is fundamentally fragile when model fidelity varies across temporal scales.
Hierarchical approaches that avoid direct dt-conditioning of the critic, or architectures that improve multi-step prediction fidelity, are more promising directions.

\bibliographystyle{plain}
\begin{thebibliography}{20}

\bibitem{sutton1999options}
R.~S. Sutton, D.~Precup, and S.~Singh.
\newblock Between MDPs and semi-MDPs: A framework for temporal abstraction in reinforcement learning.
\newblock \emph{Artificial Intelligence}, 112(1-2):181--211, 1999.

\bibitem{vezhnevets2017feudal}
A.~S. Vezhnevets, S.~Osindero, T.~Schaul, N.~Heess, M.~Jaderberg, D.~Silver, and K.~Kavukcuoglu.
\newblock FeUdal Networks for hierarchical reinforcement learning.
\newblock In \emph{ICML}, 2017.

\bibitem{hafner2023dreamerv3}
D.~Hafner, J.~Pasukonis, J.~Ba, and T.~Lillicrap.
\newblock Mastering diverse domains through world models.
\newblock \emph{arXiv preprint arXiv:2301.04104}, 2023.

\bibitem{hansen2024tdmpc2}
N.~Hansen, H.~Su, and X.~Wang.
\newblock TD-MPC2: Scalable, robust world models for continuous control.
\newblock In \emph{ICLR}, 2024.

\bibitem{nhu2025tawm}
N.~Nhu, Q.~Le, D.~Hafner, et~al.
\newblock Temporal abstraction in world models via MPC.
\newblock In \emph{ICML}, 2025.

\bibitem{janner2019mbpo}
M.~Janner, J.~Fu, M.~Zhang, and S.~Levine.
\newblock When to trust your model: Model-based policy optimization.
\newblock In \emph{NeurIPS}, 2019.

\bibitem{lu2018nondelusional}
T.~Lu, D.~Schuurmans, and C.~Boutilier.
\newblock Non-delusional Q-learning and value-iteration.
\newblock In \emph{NeurIPS}, 2018.

\bibitem{sharma2017figar}
S.~Sharma, A.~S. Lakshminarayanan, and B.~Ravindran.
\newblock Fine grained action repetition for deep reinforcement learning.
\newblock In \emph{ICLR Workshop}, 2017.

\bibitem{kim2023thick}
S.~Kim, S.~Oh, and S.~Yun.
\newblock THICK: Temporal hierarchies in world models.
\newblock \emph{arXiv preprint}, 2023.

\bibitem{deng2024hieros}
P.~Deng, A.~Patterson, and M.~White.
\newblock Hieros: Hierarchical imagination for model-based RL.
\newblock In \emph{ICML}, 2024.

\bibitem{henderson2018matters}
P.~Henderson, R.~Islam, P.~Bachman, J.~Pineau, D.~Precup, and D.~Meger.
\newblock Deep reinforcement learning that matters.
\newblock In \emph{AAAI}, 2018.

\bibitem{bjorck2022instability}
J.~Bjorck, C.~Gomes, and K.~Weinberger.
\newblock Is high variance inevitable in RL? A case study in continuous control.
\newblock In \emph{ICLR}, 2022.

\bibitem{agarwal2021deep}
R.~Agarwal, M.~Schwarzer, P.~S. Castro, A.~C. Courville, and M.~Bellemare.
\newblock Deep reinforcement learning at the edge of the statistical precipice.
\newblock In \emph{NeurIPS}, 2021.

\end{thebibliography}

\end{document}
